\subsection{1. Kết quả đạt được}
\begin{itemize}[label={--}, itemsep=1pt]
    \item \textbf{Xây dựng thành công kiến trúc đồ thị đa phương thức:} Đề tài đã đề xuất và cài đặt hoàn chỉnh mô hình học sâu kết hợp giữa MGAT và mạng Kolmogorov-Arnold (FastKAN). Kiến trúc này cho phép tích hợp hiệu quả ba phương thức (hình ảnh, văn bản, âm thanh) vào một không gian biểu diễn thống nhất, khắc phục được hạn chế về nhập nhằng ngữ cảnh của các phương pháp chỉ dựa trên hình ảnh truyền thống.
    
    \item \textbf{Cải thiện đáng kể kết quả phân lớp:} Kết quả thực nghiệm chứng minh mô hình đề xuất cho ra kết quả vượt trội hơn so với các mô hình CNN Baseline đơn phương thức, đặc biệt ở các lớp dữ liệu có độ nhập nhằng cao. Việc bổ sung thông tin từ văn bản và âm thanh đã giúp hệ thống đạt độ chính xác lên tới 94.35\% với DenseNet121 + MGAT + FastKAN, cao hơn 4\% so với MobileNetV3 Large baseline (với độ chính xác 90.28\%) khẳng định tính đúng đắn và hiệu quả của hướng tiếp cận đa phương thức.
    
    \item \textbf{Triển khai ứng dụng thực tiễn:} Đề tài đã hoàn thiện ứng dụng web ICH Predictor với giao diện trực quan, cho phép người dùng tương tác linh hoạt với hệ thống (tải lên dữ liệu đa dạng, tùy chọn mô hình). Đặc biệt, tính năng trực quan hóa cấu trúc đồ thị giúp tăng tính minh bạch trong quá trình ra quyết định của mô hình, mang lại giá trị tham khảo cho người dùng.
\end{itemize}